% Section 1.4, from lectures/L01/appendix/c_exercises.md.

\section{Exercises}
\label{c1:sec:exercises}\label{c1:app:c}

\begin{aside}
\textbf{How to work these.} Exercise~\ref{c1:ex:1} reinforces \secref{c1:sec:uart_def};
Exercises~\ref{c1:ex:2} and~\ref{c1:ex:3} go with \secref{c1:sec:uart_top}. \code{uart_def} is
\textbf{provided} and only needs reading. \code{uart_top} is yours, and Exercise~\ref{c1:ex:2} is
self-contained: everything needed to type \file{hw/uart_top.vhd}, from the entity down to the last
internal signal, is below. \secref{c1:sec:uart_top} is where the reasoning behind it lives, and
Exercise~\ref{c1:ex:3} asks for that reasoning back.

\textbf{How to check your work.} Neither module has a testbench of its own in this chapter.
\code{uart_def} is a package of constants, and \code{uart_top}'s system testbench,
\code{uart_top_tb}, stays skipped until its datapath exists in \chapref{5}. So the check here is
that \code{uart_top} \textbf{analyzes cleanly} against the transport it instantiates. Run GHDL from
\file{hw/}; the preface's \emph{Setting up} has the full flow and what \code{--assert-level=error}
does.

The answers to every part that is not code are in \secref{sol:c1}. Try each part before you read its
answer.
\end{aside}

% ------------------------------------------------------------------------------------------
\exercise{\code{uart_def}}{Reasoning}
\label{c1:ex:1}

\xpart{a} \code{uart_def.vhd} is provided. Read it against \secref{c1:sec:uart_def}, then, without
looking, name the seven register indices and the \code{STATUS}, \code{CTRL} and \code{ERROR_FLAGS}
bit positions. Analyze it, so that the names are in your \code{work} library for everything that
follows:

\begin{shell}
cd hw
ghdl -a --std=93 uart_def.vhd
\end{shell}

\noindent There is nothing to elaborate or run; a clean analysis means the package compiles and its
names are available to every module that later writes \code{use work.uart_def.all}.

\xpart{b} These are bit \textbf{positions}, not masks: \code{ST_RX_VALID} is \code{1}, not
\code{2}. Show how you would test the RX-valid bit of a status vector \code{status} using the
constant, and contrast it with how the C++ driver forms the same test
(\code{1U << status::RX_VALID}). Why does storing positions keep both sides identical to the
protocol's ``bit N'' wording?

\xpart{c} A teammate swaps two index values, setting \code{REG_RX_POP} to 6 and
\code{REG_ERR_FLAGS} to 5. Nothing in this chapter notices. At what point in the course would the
swap first surface as a failure, and in which testbench, if any? Say why it does not surface
before then.

% ------------------------------------------------------------------------------------------
\exercise{\code{uart_top}}{VHDL}
\label{c1:ex:2}
Build the skeleton in the order the file is written: the entity, the declarations, then the three
instantiations and the one placeholder assignment. Each part below carries what that step needs. Do
not instantiate \code{baud_gen}, \code{uart_tx}, \code{sync}, \code{uart_rx} or \code{uart_regs}
yet, and do not declare their signals: those blocks arrive in Chapters~2 to~5, and each brings its
own wiring with it.

\xpart{a}\parttitle{The entity.} Declare the eight ports in exactly this order. \code{uart_top_tb}
binds them by position, so the numbers are the contract; the names are yours.

\begin{narrowtable}{@{}rlll@{}}
  \thead{\#} & \thead{Port} & \thead{Dir} & \thead{Type} \\
  \midrule
  1 & \code{clock} & in & \code{std_logic} \\
  2 & \code{reset_n} & in & \code{std_logic} \\
  3 & \code{sclk} & in & \code{std_logic} \\
  4 & \code{mosi} & in & \code{std_logic} \\
  5 & \code{ss} & in & \code{std_logic} \\
  6 & \code{rx} & in & \code{std_logic} \\
  7 & \code{miso} & out & \code{std_logic} \\
  8 & \code{tx} & out & \code{std_logic} \\
\end{narrowtable}

The file needs \code{ieee.std_logic_1164} for \code{std_logic}. It does not need
\code{ieee.numeric_std} yet, since nothing here converts a vector to an integer, and it does
\textbf{not} need \code{use work.uart_def.all}: the top wires the register bus but never names a bit
on it, and \code{uart_regs} (\chapref{5}) is the only module that does. Follow the provided files'
house style: a banner comment listing the inputs and outputs, then
\code{architecture behaviour of uart_top is}.

\xpart{b}\parttitle{The internal signals.} Declare only the signals \textbf{this} chapter wires.
The method matters more than the list: open the entity of every module you are about to
instantiate, and for each port that does not connect straight to a \code{uart_top} port, declare a
signal of the same type to carry it. Here that means \code{reset_sync}, \code{spi_slave} and
\code{spi_reg_bridge}, and nothing else. Each later chapter adds its own signals the same way, when
it adds the block that needs them.

Two naming conventions run through the whole design, and both are worth adopting now. Signals that
carry SPI traffic between the two transport blocks take a \textbf{\code{spi_} prefix}. Signals that
have been through a two-flop synchronizer take an \textbf{\code{_s2} suffix}, which is why the
provided modules all expect \code{reset_s2_n} rather than \code{reset_n}.

\begin{booktable}{@{}p{6.8em}p{15.2em}p{7.6em}L@{}}
  \thead{Signal} & \thead{Type} & \thead{Driven by} & \thead{Read by} \\
  \midrule
  \code{reset_s2_n} & \code{std_logic} & \code{reset_sync}, in c) & every submodule \\
  \code{spi_rx_data} & \code{std_logic_vector(7 downto 0)} & \code{spi_slave} &
    \code{spi_reg_bridge} \\
  \code{spi_rx_valid} & \code{std_logic} & \code{spi_slave} & \code{spi_reg_bridge} \\
  \code{spi_ss_active} & \code{std_logic} & \code{spi_slave} & \code{spi_reg_bridge} \\
  \code{spi_tx_data} & \code{std_logic_vector(7 downto 0)} & \code{spi_reg_bridge} &
    \code{spi_slave} \\
  \code{reg_addr} & \code{std_logic_vector(3 downto 0)} & \code{spi_reg_bridge} &
    \code{uart_regs} (\chapref{5}) \\
  \code{reg_wdata} & \code{std_logic_vector(31 downto 0)} & \code{spi_reg_bridge} &
    \code{uart_regs} (\chapref{5}) \\
  \code{reg_write} & \code{std_logic} & \code{spi_reg_bridge} & \code{uart_regs} (\chapref{5}) \\
  \code{reg_rdata} & \code{std_logic_vector(31 downto 0)} & zeros, in d); then \code{uart_regs}
    (\chapref{5}) & \code{spi_reg_bridge} \\
\end{booktable}

The four register-bus signals are here even though the block that answers on them is four chapters
away, because \code{spi_reg_bridge} has ports for them and an instantiation must connect every port.
That is the register bus existing with nothing behind it, which is the shape this chapter is really
about.

\code{reg_addr} is four bits, not three, even though there are only seven registers: the protocol's
command byte reserves bits 3 to 0 for the index.

\xpart{c}\parttitle{The reset synchronizer.} \code{reset_n} arrives from off-chip, so it must
assert asynchronously and release synchronously, two flip-flops later. That job is done by the
provided \code{reset_sync}, which you instantiate rather than write; its two internal flops are its
own business, so the only signal it adds to your top is \code{reset_s2_n}. Everything inside
\code{uart_top} takes \code{reset_s2_n}, and nothing but this instance sees \code{reset_n} directly.

\bookfigure{reset_sync}{Module \code{reset_sync}}{c1:fig:reset_sync}

\begin{narrowtable}{@{}rlll@{}}
  \thead{\#} & \thead{\code{reset_sync} port} & \thead{Dir} & \thead{Connect to} \\
  \midrule
  1 & \code{clock} & in & \code{clock} \\
  2 & \code{reset_n} & in & \code{reset_n} \\
  3 & \code{reset_s2_n} & out & \code{reset_s2_n} \\
\end{narrowtable}

\begin{vhdlcode}
reset_sync: entity work.reset_sync
    port map(clock, reset_n, reset_s2_n);
\end{vhdlcode}

\xpart{d}\parttitle{The transport.} Instantiate \code{spi_slave} and \code{spi_reg_bridge}, both
positionally. The last column names a signal from part b) or a port from part a).

\bookfigure{spi_slave}{Module \code{spi_slave}}{c1:fig:spi_slave}

\begin{booktable}{@{}rp{7.4em}p{1.8em}p{15.2em}L@{}}
  \thead{\#} & \thead{\code{spi_slave} port} & \thead{Dir} & \thead{Type} & \thead{Connect to} \\
  \midrule
  1 & \code{clock} & in & \code{std_logic} & \code{clock} \\
  2 & \code{reset_s2_n} & in & \code{std_logic} & \code{reset_s2_n} \\
  3 & \code{sclk} & in & \code{std_logic} & \code{sclk} \\
  4 & \code{mosi} & in & \code{std_logic} & \code{mosi} \\
  5 & \code{ss} & in & \code{std_logic} & \code{ss} \\
  6 & \code{tx_data} & in & \code{std_logic_vector(7 downto 0)} & \code{spi_tx_data} \\
  7 & \code{miso} & out & \code{std_logic} & \code{miso} \\
  8 & \code{rx_data} & out & \code{std_logic_vector(7 downto 0)} & \code{spi_rx_data} \\
  9 & \code{rx_valid} & out & \code{std_logic} & \code{spi_rx_valid} \\
  10 & \code{ss_active} & out & \code{std_logic} & \code{spi_ss_active} \\
\end{booktable}

\bookfigure{spi_reg_bridge}{Module \code{spi_reg_bridge}}{c1:fig:spi_reg_bridge}

\begin{booktable}{@{}rp{9.8em}p{1.8em}p{15.2em}L@{}}
  \thead{\#} & \thead{\code{spi_reg_bridge} port} & \thead{Dir} & \thead{Type} &
    \thead{Connect to} \\
  \midrule
  1 & \code{clock} & in & \code{std_logic} & \code{clock} \\
  2 & \code{reset_s2_n} & in & \code{std_logic} & \code{reset_s2_n} \\
  3 & \code{ss_active} & in & \code{std_logic} & \code{spi_ss_active} \\
  4 & \code{rx_data} & in & \code{std_logic_vector(7 downto 0)} & \code{spi_rx_data} \\
  5 & \code{rx_valid} & in & \code{std_logic} & \code{spi_rx_valid} \\
  6 & \code{reg_rdata} & in & \code{std_logic_vector(31 downto 0)} & \code{reg_rdata} \\
  7 & \code{tx_data} & out & \code{std_logic_vector(7 downto 0)} & \code{spi_tx_data} \\
  8 & \code{reg_addr} & out & \code{std_logic_vector(3 downto 0)} & \code{reg_addr} \\
  9 & \code{reg_wdata} & out & \code{std_logic_vector(31 downto 0)} & \code{reg_wdata} \\
  10 & \code{reg_write} & out & \code{std_logic} & \code{reg_write} \\
\end{booktable}

The names that look shared are not one bus. \code{spi_slave}'s \code{rx_data} is an output and the
bridge's \code{rx_data} is an input, so \code{spi_rx_data} runs from the slave to the bridge;
\code{spi_tx_data} runs the other way, from the bridge's \code{tx_data} output into the slave's
\code{tx_data} input. Read them as two one-directional lanes and the wiring falls out.

Nothing drives \code{reg_rdata} yet, so give it a value:

\begin{vhdlcode}
reg_rdata <= (others => '0');
\end{vhdlcode}

\noindent The file would analyze without this; an undriven signal is legal. It is here so that the
bridge reads a defined value instead of \code{'U'} while the bank is missing. \textbf{Delete this
line in \chapref{5}}, when \code{uart_regs} starts driving \code{reg_rdata}. Leave it in and the
vector has two drivers: every \code{'1'} from the bank resolves against the \code{'0'} here to
\code{'X'}, and the system testbench fails on a read-back caused by a line you wrote four chapters
earlier.

\xpart{e}\parttitle{The check.} Analyze the skeleton against the package and the transport it
instantiates:

\begin{shell}
cd hw
ghdl -a --std=93 uart_def.vhd reset_sync.vhd spi_slave.vhd spi_reg_bridge.vhd uart_top.vhd
\end{shell}

\noindent Then run \code{make build-vhdl} from the repository root. Confirm that \code{uart_def} and
\code{uart_top} both analyze cleanly, that the two transport testbenches pass (they are the only
ones that can run this early, since they check provided modules), and that \code{uart_top_tb}
reports \textbf{skipped}. Which modules does the skip message name as missing, and which chapter
(2, 3, 4 or 5) builds each? One of the six is never instantiated by \code{uart_top} at all; find
it, and say why the build needs it anyway.

% ------------------------------------------------------------------------------------------
\exercise{What a clean analysis does not prove}{Reasoning}
\label{c1:ex:3}
\code{uart_top} analyzes, and almost nothing about it is proven. That gap is the lesson of this
chapter.

\xpart{a} Swap two same-type ports in your entity, say \code{sclk} and \code{mosi}, and re-analyze.
Does \code{ghdl -a} catch it? If not, at which point in the course does the mistake finally show up,
and what does that tell you about the cost of a positional contract?

\xpart{b} \code{reg_rdata} is tied to zeros and no register bank answers on the bus. What does an
SPI read transaction return with that placeholder in place, and why is that harmless until
\chapref{5}? What would it return instead if you left the placeholder out altogether?

\xpart{c} The provided \code{reset_sync} asserts asynchronously but releases synchronously. Suppose
it released asynchronously too, wiring \code{reset_n} straight to every block. Describe the failure
that a near-edge release could cause. Then explain why \code{sync} (\chapref{3}) cannot do this job,
given that it is also two flip-flops in series: look at its port list and say what it would need
that only \code{reset_sync} can give it.

\xpart{d} Why does building \code{uart_top} before any block it instantiates fix each block's port
list in advance? What has to be true about \code{baud_gen}'s ports in \chapref{2} for its
instantiation here to bind without you touching the top again?
